% ============================================================================
\chapter{Entity-Component System (ECS)}
% ============================================================================

\section{Design Philosophy}

The ECS in Caffeine is \emph{archetype-based}: entities with the same
set of component types share a single \texttt{Archetype}. This groups
the data for all entities with a given composition contiguously in memory,
enabling iteration over a specific component type to proceed without cache
misses.

\begin{definition}[Entity]
An entity is an unsigned 32-bit integer identifier. It carries no data.
All observable state is stored in components associated with the entity.
An entity with index $\texttt{u32\_max}$ is conventionally \emph{invalid}.
\end{definition}

\begin{definition}[Component]
A component is a plain-old-data (POD) struct. It must have no virtual
methods and no hidden heap allocation. The engine enforces this by
instantiating components through typed \texttt{ComponentPool<T>} which
stores them in a \texttt{Vector<T>}.
\end{definition}

\begin{definition}[Archetype]
An archetype $\mathcal{A}$ is the set of component types associated with a
group of entities. Two entities belong to the same archetype if and only if
they possess exactly the same set of component types.
\end{definition}

\section{Component Identification}

Each component type receives a unique 32-bit ID at program initialisation
via \texttt{ComponentID::get<T>()}:

\begin{lstlisting}[caption={Type-ID registration}]
template<typename T>
static u32 get() {
    static const u32 id = nextID();  // initialised once per type
    return id;
}
\end{lstlisting}

The \texttt{static} local variable is guaranteed to be initialised exactly
once (C++11 magic statics). The counter is an \texttt{std::atomic<u32>}
to support concurrent registration from multiple translation units.
Up to 256 distinct component types are permitted.

\section{Component Set (\texttt{ComponentSet})}

A \texttt{ComponentSet} is a bitmask of component IDs. Archetype identity is
determined by comparing two \texttt{ComponentSet}s. Set operations (union,
intersection, subset test) reduce to bitwise OR, AND, and masking.

\section{Archetype Internal Structure}

Each \texttt{Archetype} contains:
\begin{itemize}
  \item A \texttt{Vector<u32>} of entity IDs (the entities belonging to it).
  \item A \texttt{PoolRegistry}: a fixed-size array of 64 pool entries,
        indexed by component ID, each entry holding a \texttt{void*} pool
        pointer and function pointers for copy, create, and remove.
\end{itemize}

\subsubsection{Removal by Swap-and-Pop}

When an entity at index $i$ is removed from an archetype, the entity at
the last position $n-1$ is moved into slot $i$. This preserves contiguity
in $O(1)$ without shifting all subsequent elements.

\begin{equation}
\text{entity}[i] \leftarrow \text{entity}[n-1]; \quad n \leftarrow n-1
\end{equation}

The same swap-and-pop is applied to every component pool for that
archetype.

\section{Command Buffer}

Structural mutations (creating or destroying entities, adding or removing
components) \emph{during iteration} would invalidate in-flight iterators.
The \texttt{CommandBuffer} defers all mutations by recording them as
\texttt{ICommand} objects:

\begin{lstlisting}[caption={Deferred entity creation}]
Entity CommandBuffer::create() {
    u32 pendingID = m_nextPendingID++;
    m_commands.pushBack(make_unique<CreateCommand>(pendingID));
    return Entity(pendingID, nullptr);  // not yet alive in world
}
\end{lstlisting}

\texttt{execute(World\&)} is called at a safe point (end of frame,
after all systems finish) to replay all recorded mutations on the live world.

\begin{specbox}{Invariant 6.1 --- Structural mutation safety}
  ECS structural mutations (entity creation/destruction, component
  add/remove) during iteration \textbf{must} go through
  \texttt{CommandBuffer}. Direct world mutation during a
  \texttt{forEach} loop is undefined behaviour.
\end{specbox}

\section{Systems}

Systems implement \texttt{ISystem} and are registered in
\texttt{SystemRegistry}. Each system exposes \texttt{onUpdate(World\&, f32 dt)}.
The world provides \texttt{forEach<C1, C2, \ldots>} which iterates over all
entities possessing the queried component types, invoking a callable
with typed references.
